\documentclass[../../../include/open-logic-section]{subfiles}

\begin{document}

\olfileid[hi]{his}{set}{infinitesimals}
\olsection{अनंतसूक्ष्म राशियाँ और अवकलन}

न्यूटन और लाइबनिट्स ने सत्रहवीं शताब्दी के अंत में (स्वतंत्र रूप से) कलन की
खोज की। कलन का एक विशेष महत्त्वपूर्ण अनुप्रयोग \emph{अवकलन} था। मोटे तौर
पर, अवकलन का उद्देश्य किसी बिंदु पर फलन की ``परिवर्तन दर'', या प्रवणता, की
संकल्पना देना है।

इस विचार को स्पष्ट करने का एक सजीव ढंग यह है। नीचे काले रंग में दिखाए गए फलन
$f(x) = \nicefrac{x^2}{4} + \nicefrac{1}{2}$ पर विचार करें:
\begin{center}
	\begin{tikzpicture}[scale=1]
	\draw[->, gray] (-1,0) -- (4.25,0) node[right] {$x$};
	\draw[->, gray] (0,-0.25) -- (0,5.25) node[above] {$f(x)$};
	
	\foreach \x/\xtext in {1/1, 2/2, 3/3, 4/4}
	\draw[shift={(\x,0)}, gray] (0pt,2pt) -- (0pt,-2pt) node[below] {$\xtext$};
	
	\foreach \y/\ytext in {1/1, 2/2, 3/3, 4/4, 5/5}
	\draw[shift={(0,\y)}, gray] (2pt,0pt) -- (-2pt,0pt) node[left] {$\ytext$};
	
	\draw[oldiagcolorC] (0.5, 9/16) -- (3.5, 57/16) -- (3.5, 9/16)--cycle;
	\draw[oldiagcolorD] (.5, 9/16) -- (2.5, 33/16) -- (2.5, 9/16) -- cycle;
	\draw[oldiagcolorE] (.5, 9/16) -- (1.5, 17/16) -- (1.5, 9/16) -- cycle;
	\draw[thick] (-1,.75) parabola bend (0,0.5) (4,4.5);
	\end{tikzpicture}
\end{center}
मान लें हम $c = \nicefrac{1}{2}$ पर फलन की प्रवणता ज्ञात करना चाहते हैं। हम
ऐसा त्रिभुज खींचकर आरंभ करते हैं जिसका कर्ण उस बिंदु पर प्रवणता का सन्निकटन
करता है, संभवतः ऊपर का लाल त्रिभुज। जब $\beta$ हमारे त्रिभुज की आधार-लंबाई
है, तो उसकी ऊँचाई $f(\nicefrac{1}{2}+\beta) - f(\nicefrac{1}{2})$ है, अतः
कर्ण की प्रवणता है:
\[
\frac{f(\nicefrac{1}{2}+\beta) - f(\nicefrac{1}{2})}{\beta}.
\]
अतः आधार-लंबाई~$3$ वाले हमारे !!{colorC} त्रिभुज की प्रवणता ठीक~$1$ है।
छोटे त्रिभुज, आधार-लंबाई~$2$ वाले !!{colorD} त्रिभुज, का कर्ण बेहतर
सन्निकटन देता है; उसकी प्रवणता $\nicefrac{3}{4}$ है। उससे भी छोटा,
आधार-लंबाई~$1$ वाला !!{colorE} त्रिभुज और बेहतर सन्निकटन देता है; उसकी
प्रवणता $\nicefrac{1}{2}$ है।

लगातार छोटे होते त्रिभुज हमें लगातार बेहतर सन्निकटन देते हैं। इसलिए हम कुछ
ऐसा कह सकते हैं: \emph{अनंतसूक्ष्म} आधार-लंबाई वाले त्रिभुज का कर्ण हमें $c
= \nicefrac{1}{2}$ पर स्वयं प्रवणता देता है। इस प्रकार हमें बिंदु $c$ पर
फलन $f$ के (प्रथम) अवकलज का सूत्र मिलेगा:
\[
{f'}(c) = \frac{f(c+\beta) - f(c)}{\beta} \text{ जहाँ $\beta$ अनंतसूक्ष्म है।}
\]
और मोटे तौर पर न्यूटन और लाइबनिट्स ने यही कहा था।

किंतु उन्होंने यह कहा है, तो हमें उनसे पूछना होगा: \emph{अनंतसूक्ष्म} क्या
है? एक गंभीर दुविधा उत्पन्न होती है। यदि $\beta = 0$, तो $f'$ अपरिभाषित है,
क्योंकि इसमें $0$ से भाग देना पड़ता है। किंतु यदि $\beta > 0$, तो हमें केवल
प्रवणता का \emph{सन्निकटन} मिलता है, स्वयं प्रवणता नहीं।

यह कोई काल-विसंगत चिंता नहीं है। यहाँ बर्कली न्यूटन के अनुयायियों की आलोचना
करते हैं:
\begin{quote}
मैं मानता हूँ कि चिह्नों से किसी वस्तु अथवा कुछ भी नहीं, दोनों में से किसी
को निरूपित कराया जा सकता है; और फलतः मूल संकेतन $c + \beta$ में $\beta$
किसी वृद्धि अथवा कुछ भी नहीं, दोनों में से किसी का अर्थ दे सकता था। किंतु
फिर आप उससे इनमें से जिसका भी अर्थ दिलाएँ, आपको उसी अर्थ के अनुरूप संगत तर्क
करना चाहिए और दोहरे अर्थ के आधार पर आगे नहीं बढ़ना चाहिए: ऐसा करना स्पष्ट
कुतर्क होगा। (\citealt[\S{}XIII]{Berkeley1734}, चरों को पिछले पाठ के अनुरूप
बदला गया है)
\end{quote}
बर्कली के विरुद्ध अनंतसूक्ष्म कलन का बचाव करने के लिए कोई उत्तर दे सकता है
कि ``अनंतसूक्ष्म राशियों'' की बात केवल आलंकारिक है। कहा जा सकता है कि जब तक
हम \emph{वास्तव में छोटा} त्रिभुज लेते हैं, हमें स्पर्शरेखा का \emph{पर्याप्त
अच्छा} सन्निकटन मिलेगा। बर्कली के पास इसका भी उत्तर था: यद्यपि वह
अभियांत्रिकी के लिए पर्याप्त अच्छा हो सकता है, पर वह गणित की \emph{स्थिति}
को दुर्बल करता है, क्योंकि
\begin{quote}
हमें बताया जाता है कि \emph{in rebus mathematicis errores qu\`{a}m minimi
non sunt contemnendi}। [गणित के मामले में छोटी-से-छोटी त्रुटियों की उपेक्षा
नहीं की जानी चाहिए।] \citep[\S{}IX]{Berkeley1734}
\end{quote}
तिरछे अक्षरों वाला अंश स्वयं न्यूटन की \emph{Quadrature of Curves} (1704)
से लगभग शब्दशः उद्धरण है।

बर्कली की दार्शनिक आपत्तियाँ अत्यंत मर्मभेदी हैं। फिर भी कलन एक अत्यधिक सफल
उद्यम था और गणितज्ञ बिना त्रुटि में पड़े उसका उपयोग करते रहे।

\end{document}
